To support application operators in detecting failures and identiying their possible root causes, we surveyed existing techniques for anomaly detection and root cause analysis  in modern multi-service applications. 
Other than presenting the methods applied by such techniques, we discussed the type and granuality of anomalies they can detect/analyse, their setup costs (\viz the instrumentation that application operators must apply to their applications to use such techniques, and the additional artifacts they must produce), their accuracy and applicability to dynamically changing applications, and open challenges on the topic.

We believe that our survey can provide benefits to practicitioners working with modern multi-service applications, such as microservice-based applications, for instance. 
Our survey provides a first support for setting up a pipeline of anomaly detection and root cause analysis techniques, so that failures are automatically detected and analysed in multi-service applications.
We indeed not only highlighted which existing techniques already come in an integrated pipeline, but we also discussed the type and granularity of anomalies they consider. 
\upd{This provides a baseline for application operators to identify the anomaly detection and root cause analysis techniques that ---even if coming as independent solutions--- could be integrated in a pipeline, since they identify and explain the same type/granularity of anomalies}.
This, togeter with our discussion on the techniques' setup costs (\Cref{sec:detection:discussion:type-granularity-costs,sec:rca:discussion:root-causes-costs}), provides a first support to practitioners for choosing the most suited techniques for their needs.

In this perspective, a quantitative comparison of the performance of the surveyed techniques would further support practitioners, and it definitely deserves future work. 
Indeed, whilst the surveyed techniques already measure their performance or accuracy, sometimes also in comparison with other existing techniques, this is typically done with experiments on different reference applications running in different environments.
This makes it complex to quantitatively compare the existing anomaly detection and root cause analysis techniques, hence opening the need for studies quantitatively comparing the performance of such techniques by fixing the reference applications and runtime environments, much in a similar way as Arya \etal \cite{Arya2021_EvaluationCausalInferenceAIOps} did for a subset of existing root cause analysis techniques.  
Quantitative performance comparisons would complement the results of our qualitative comparison, providing practitioners with more tools to choose the anomaly detection and root cause analysis techniques most suited to their applications' requirements.

Our survey can also help researchers investigating failures and their handling in modern multi-service applications, as it provides a first structured discussion of anomaly detection and root cause analysis techniques in such a kind of applications.
For instance, we discussed how false positives and false negatives can negatively affect the accuracy of observed anomalies/identified root causes (\Cref{sec:detection:discussion:adaptability,sec:rca:discussion:accuracy}).
Whilst false positives could result in unnecessary work for application operators, false negatives could severely impact on the outcomes of the surveyed technique: false negatives indeed correspond to considering a service as not anomalous or not causing an observed anomaly, even if this was actually the case.
In both cases, false positives/negatives can occur if the services forming an application or its runtime conditions change over time.
Whilst some of the surveyed techniques already consider the issue of accuracy losses in the case of application changing over time, this is done with time consuming processes, \eg by re-training the machine learning models used to detect anomalies \cite{14_Samir2019_DLA,26_Wang2020_WorkflowAwareFaultDiagnosis,17_Gan2019_Seer} or by relying on application operators to provide updated artifacts to identify possible root causes \cite{11_Brandon2020_GraphBasedRootCauseAnalysis}.
Re-training or manual updates, as well as the other proposed updates, are hence not suited to get continuously enacted, but rather require to find a suitable tradeoff between the update period and the accuracy of enacted anomaly detection/root cause analysis \cite{14_Samir2019_DLA,17_Gan2019_Seer}.
However, modern software delivery practices are such that new versions of the services forming an application are continuously released \cite{Humble2010_ContinuousDelivery}, and cloud application deployments are such that services are often migrated from one runtime to another \cite{Carrasco2020_MigrationTransCloud}.  
An interesting research direction is hence to devise anomaly detection and root cause analysis techniques that can effectively work also in presence of continuous changes in multi-service applications, either adapting the surveyed ones or by devising new ones, \eg by exploiting the recently proposed continual learning solutions \cite{Lomonaco2021_ContinualAI}.

Other two interesting research directions already emerged in \Cref{sec:detection:discussion:explainability-countermeasures,sec:rca:discussion:explainability-countermeasures}, where we discussed the need for explainability and countermeasures.
On the one hand, associating observed anomalies and their identified root causes with explanations of why they are considered so would allow application operators to directly exclude false positives, which in turn means that, \eg we could avoid cutting the possible root causes for an observed anomaly to only the most probable ones.
Indeed, an application operator could be provided with all possible root causes and she could later decide which deserve to be troubleshooted based on their explanations.
This hence calls for anomaly detection and root cause analysis techniques that are \enquote{explainable by design}, much in the same way as the need for explainability is nowadays recognized in AI \cite{Guidotti2019_Explainability}.

On the other hand, recommending potential countermeasures to be enacted to avoid (the failures corresponding) to observed anomalies and their root causes would enable avoiding such failures to happen again in the future.
This could be done both for anomaly detection and for root cause analysis.
Indeed, whilst avoiding (the failure corresponding to a) detected anomaly of course needs acting on the root causes for such an anomaly, countermeasures could anyhow be taken to avoid the anomaly in a service to propagate to other services.
At the same time, the possible root causes for an observed anomaly are identified based on data collected on the services forming an application, which is processed only to enact root cause analysis.
Such data could however be further processed to automatically extract possible countermeasures to avoid it to cause the observed anomaly.
For instance, if a service logged an error event that caused a cascade of errors reaching the service where the anomaly was observed, the root cause analysis technique could further process the error logs to suggest how to avoid such error propagation to happen again, \eg by introducing circuit breakers or bulkheads \cite{Richardson2018_MicroservicesPatterns}.
If a service was responding too slowly and caused some performance anomaly in another service, this information can be used to train machine learning models to predict similar performance degradations in the root causing service, which could then be used to preemptively scale such service and avoid the performance degradation to again propagate in future runs of an application.
The above are just two out of many possible examples, deserving further investigation and opening a new research direction. 

% Finally, it is worth mentioning an alternative, which ---by changing the perspective for enacting failure detection and root cause analysis--- could result in explainable-by-design techniques that could possibly avoid false positives/negatives to occur.
% All the surveyed techniques follow the \enquote{medical} metaphore: the occurrence of a failure in a service is detected based on whether the service is experiencing some symptoms, \viz anomalies.
% Its propagation from/to other services is analysed by checking whether such services are also experiencing similar symptoms, \viz anomalies.
% Another perspective could be to devise methodologies where both functional and performance failures are, \eg explicitly logged by the services themselves, possibly also indicating their severity in accordance with the Syslog standard \cite{Syslog2009}.
% Failures (rather than their symptoms) would be directly observable in service logs, which could be automatically analysed to determine causal relationships among failures, hence determining the possible root causes for a considered failure. 
% We are aware that this is quite far from how things are currently done, as, \eg while the surveyed techniques rely on available application logs or on common instrumentation for modern multi-service application, such an approach would require to change the service in an application log events.
% We anyhow believe it deserves further investigation, as it could change and enhance the way we currently deal with failures in modern multi-service applications.